\chapter{Code Appendix}
%
\label{chap:CodeAppendix}
%
This appendix contains key code implementations from the thesis project. The complete source code is available in the project repository.

\section{Document Conversion Module}
%
\subsection{PDF Text Extraction}
%
The following code demonstrates PDF text extraction using PyMuPDF:

\begin{verbatim}
import fitz  # PyMuPDF
import os

def extract_pdf_text(pdf_path):
    """Extract text from PDF using PyMuPDF"""
    text = ""
    try:
        doc = fitz.open(pdf_path)
        for page in doc:
            text += page.get_text()
        doc.close()
    except Exception as e:
        print(f"Error processing {pdf_path}: {e}")
    return text

# Usage
pdf_text = extract_pdf_text("document.pdf")
\end{verbatim}

\section{Language Detection Module}
%
\subsection{Langdetect Integration}
%
Language detection using the langdetect library:

\begin{verbatim}
from langdetect import detect, LangDetectException

def detect_language(text):
    """Detect language of text"""
    try:
        language = detect(text)
        return language
    except LangDetectException:
        return "unknown"

# Usage
text = "Dies ist ein deutscher Text."
lang = detect_language(text)
print(f"Detected language: {lang}")
\end{verbatim}

\section{Database Operations}
%
\subsection{SQLite Database Schema}
%
Creating and populating the document database:

\begin{verbatim}
import sqlite3

def create_document_database(db_path):
    """Create database for document metadata"""
    conn = sqlite3.connect(db_path)
    cursor = conn.cursor()
    
    cursor.execute('''
        CREATE TABLE IF NOT EXISTS documents (
            id INTEGER PRIMARY KEY,
            file_name TEXT NOT NULL,
            original_format TEXT,
            jahr TEXT,
            language TEXT,
            processing_date TIMESTAMP,
            content_length INTEGER
        )
    ''')
    
    conn.commit()
    conn.close()

# Usage
create_document_database("documents.db")
\end{verbatim}

\subsection{Data Aggregation}
%
Aggregating document statistics:

\begin{verbatim}
import sqlite3
import pandas as pd

def aggregate_statistics(db_path):
    """Aggregate document statistics by year"""
    conn = sqlite3.connect(db_path)
    
    query = '''
        SELECT jahr, original_format, COUNT(*) as count
        FROM documents
        GROUP BY jahr, original_format
        ORDER BY jahr, count DESC
    '''
    
    df = pd.read_sql_query(query, conn)
    conn.close()
    
    return df

# Usage
stats = aggregate_statistics("documents.db")
print(stats)
\end{verbatim}

\section{Visualization Module}
%
\subsection{Matplotlib Chart Generation}
%
Creating analysis visualizations:

\begin{verbatim}
import matplotlib.pyplot as plt
import pandas as pd

def plot_format_distribution(data_dict):
    """Create stacked bar chart of format distribution"""
    fig, ax = plt.subplots(figsize=(14, 7))
    
    # Prepare data
    years = sorted(data_dict.keys())
    formats = set()
    for year_data in data_dict.values():
        formats.update(year_data.keys())
    formats = sorted(list(formats))
    
    # Create stacked bars
    bottom = [0] * len(years)
    for fmt in formats:
        values = [data_dict[year].get(fmt, 0) for year in years]
        ax.bar(years, values, label=fmt, bottom=bottom)
        bottom = [bottom[i] + values[i] for i in range(len(years))]
    
    ax.set_xlabel('Year')
    ax.set_ylabel('Document Count')
    ax.set_title('Document Format Distribution by Year')
    ax.legend()
    
    plt.tight_layout()
    plt.savefig('format_distribution.png', dpi=150)
    plt.show()
\end{verbatim}

\section{Complete Processing Pipeline}
%
\subsection{End-to-End Example}
%
Complete processing example combining modules:

\begin{verbatim}
import os
from pathlib import Path

def process_document_collection(input_dir, output_dir, db_path):
    """Complete document processing pipeline"""
    
    # Create database
    create_document_database(db_path)
    conn = sqlite3.connect(db_path)
    cursor = conn.cursor()
    
    # Process all documents
    for file_path in Path(input_dir).rglob("*"):
        if file_path.is_file():
            # Extract text based on format
            if file_path.suffix.lower() == '.pdf':
                text = extract_pdf_text(str(file_path))
            else:
                # Handle other formats
                continue
            
            # Detect language
            language = detect_language(text)
            
            # Store metadata
            cursor.execute('''
                INSERT INTO documents 
                (file_name, original_format, language, processing_date)
                VALUES (?, ?, ?, datetime('now'))
            ''', (file_path.name, file_path.suffix, language))
    
    conn.commit()
    conn.close()
    
    # Generate statistics and visualizations
    stats = aggregate_statistics(db_path)
    print(f"Processed {len(stats)} documents")

# Usage
process_document_collection(
    "input_documents/",
    "output_markdown/",
    "documents.db"
)
\end{verbatim}

\section{Jupyter Notebook Integration}
%
The main analysis is performed in the \texttt{Analysis.ipynb} Jupyter notebook, which provides:

\begin{itemize}
    \item Interactive exploration of document statistics
    \item Dynamic visualization generation
    \item Memory monitoring and resource management
    \item Comprehensive error handling and logging
    \item Progress indicators for batch operations
\end{itemize}

The notebook integrates all modules and provides a user-friendly interface for analyzing document collections.

\section{Dependencies and Requirements}
%
Key Python packages used:

\begin{itemize}
    \item \texttt{pymupdf}: PDF processing
    \item \texttt{python-docx}: DOCX format handling
    \item \texttt{langdetect}: Language detection
    \item \texttt{pandas}: Data manipulation
    \item \texttt{matplotlib}: Visualization
    \item \texttt{sqlite3}: Database operations
    \item \texttt{psutil}: System monitoring
\end{itemize}

Installation via pip:

\begin{verbatim}
pip install pymupdf python-docx langdetect pandas matplotlib psutil
\end{verbatim}

\section{Configuration and Customization}
%
The pipeline is highly customizable through:

\begin{itemize}
    \item Configuration files specifying document types and formats
    \item Database schema adjustments for additional metadata
    \item Visualization customization through matplotlib parameters
    \item Filter parameters for language and format selection
    \item Process parameters for memory and performance tuning
\end{itemize}
